\documentclass[a4paper, oneside]{article}
\usepackage{graphicx}
\usepackage{amsthm}
\usepackage{amsmath}
\usepackage{amssymb}
\usepackage[a4paper,
            bindingoffset=0.2in,
            left=2cm,
            right=2cm,
            top=2cm,
            bottom=2cm,
            footskip=.25in]{geometry}
\usepackage[italian]{babel}
\usepackage{pgfplots}
\usepackage{tabularx}
\usepackage{tikz}
\usepackage{wrapfig}
\usepackage{color}
\usepackage[d]{esvect}
\definecolor{page}{rgb}{0.129,0.157,0.212}
\pagecolor{page}
\color{white}
\graphicspath{ {./images/} }
\usetikzlibrary{shapes.geometric}
\usetikzlibrary{datavisualization}
\usetikzlibrary{datavisualization.formats.functions}
\usetikzlibrary{patterns}
\pgfplotsset{width=10cm,compat=1.18}

\title{Appunti di Lab II}
\author{Tommaso Miliani}
\date{18-03-26}

\begin{document}
\newtheoremstyle{theoremEnv}
                {}          % Space above
                {}          % Space below
                {\slshape}  % Body font
                {}          % Indent amount
                {\bfseries} % Head font
                {.}         % Punctuation after head
                {\newline}  % Space after theorem head
                {}          % Theorem head spec
\theoremstyle{theoremEnv}

\newtheorem{definition}{Definizione}[section]
\newtheorem{theorem}{Teorema}[section]
\newtheorem{lemma}{Proposizione}[section]
\newtheorem{observation}{Osservazione}[section]
\newtheorem{corollary}{Corollario}[theorem]
\newtheorem{example}{Esempio}[section]
\newtheorem{remark}{Enunciato}[section]

\maketitle

\section{Generatori di corrente}
Il generatore di corrente cerca di pompare delle cariche con il medesimo 
flusso allo stesso modo indipendentemente dalla resistenza sul circuito. 
Se si ha un sistema che prende fuoco con troppa corrente, si può regolare una 
manopola per evitare che passi troppa corrente. Il generatore $\epsilon$ cerca sempre di mantenere una 
differenza di potenziale, mentre un generatore di corrente (come un fotodiodo) cerca di 
mantenere sempre costante la corrente in un circuito. Il generatore di corrente $I$
si indica come in figura. La sua caratteristica devia dalla sua situazione ideale in modo tale 
che, indipendentemente dal carico sulla corrente 
\begin{gather*}
    \begin{tikzpicture}
        \draw[->](0, 0) -- (4, 0) node[at end, below] {$i$};
        \draw[->](0, 0) -- (0, 4) node[at end, left] {$V$};
        \draw(0, 3) -- (3, 3);
        \draw[dashed](3, 0) -- (3, 3);
        \draw(3, 0) .. controls (2.7, 1) and (2.6, 2) .. (2.5, 3);
    \end{tikzpicture}
\end{gather*}
Dunque un generatore si schematizza sempre con una resistenza, la corrente disponibile 
non sarà mai $I$ ma si avranno due correnti $i_1$ che va verso il circuito, mentre $i_2$ 
va all'interno del circuito con la resistenza. Non si avrà dunque a disposizione tutta la corrente 
nel circuito ma, noto il carico sul generatore e la resistenza, le correnti si rapporteranno 
come il rapporto delle resistenze.
\begin{gather*}
    I = i_1 + i_2 = \text{const}
\end{gather*}
Allora
\begin{gather*}
    i_1 = I - \frac{\Delta V}{R} \ \Longrightarrow \ I = i_1 + \frac{\Delta V}{R}
\end{gather*}

\section{Teorema di Norton}
Preso un generatore di tensione collegato ad una resistenza $\rho$ 
e collegato ai capi $A$ e $B$. Equivalentemente si considera un generatore 
di corrente $I$ con resistenza $R$. Si vuole trovare ora le condizioni secondo le 
quali ai capi $A$ e $B$ si ha la stessa differenza di tensione ai capi $A'$ e $B'$. 
\begin{gather*}
    V_A - V_B = \epsilon - i\rho \qquad V_A' - V_B' = i_2R
\end{gather*}
Se si vuole che le differenze di tensione siano equivalenti, allora si eguagliano
\begin{gather*}
    \epsilon - i\rho = i_2R 
\end{gather*}
si ha che 
\begin{gather*}
    I = i_1 + i_2 = i_1 + \frac{\epsilon}{R} - i\frac{\rho}{R}
\end{gather*}
Allora 
\begin{gather*}
    i_1 = I - \frac{\epsilon}{R} + i\frac{\rho}{R}
\end{gather*}
Dunque per far sì che $i = i_1$, si deve avere che $I = \frac{\epsilon}{R}$ e 
inoltre che $\rho = R$. In altre parole, se si ha un generatore di corrente con queste caratteristiche, 
lo si può sostituire con un generatore di tensione $\frac{\epsilon}{R}$ e resistenza $R$. 

\section{Amplificatore operazionale}
Nell'amplificatore operazionale si ha  circuito con 
un generatore $v_0$ pilotato e un generatore $v_s$. 
Per l'amplificatore 
\begin{gather*}
    v_0 = A(v_+ - v_-)
\end{gather*}
Si ha che
\begin{gather*}
    A_f = - \frac{R_f}{R_s} \\
    v_- \to 0 \qquad i_s = i_f
\end{gather*}
CHe si comporta come una massas virtuale. 
\begin{gather*}
    v_0 = A(v_+ - v_-) = -Av_- \\
    v_s = i_s R_s - iR_i \qquad \text{Seconda legge sulla M1}\\
    i_s + i = i_f \qquad \text{Prima legge sul nodo N}\\
    v_- = -iR_i \qquad \text{Legge di Ohm}\\
    -v_0 = iR_i + i_fR_f + i_fR_0 \qquad \text{Seconda legge su maglia M2}
\end{gather*}
Svolgendo i conti 
\begin{gather*}
    v_0 = - \frac{R_i(AR_f - R_0)}{(R_i + R_s)(R_f + R_0) + R_iR_s(i + A)}v_s
\end{gather*}
Quindi 
\begin{gather*}
    i = - v_s \frac{R_f + R_0}{(R_i + R_s) (R_f + R_0) + R_iR_s(i + A)} \to 0
\end{gather*}
Questo amplificatore, nel limite nel quale $A$ tende all'infinito, si comporta come un generatore
ideale secondo il seguente circuito:

Se si mandano tre differenze di tensione ad un amplificatore secondo il seguente disegno:

SI ha che
\begin{gather*}
    i_f = \sum_k i_k = \sum_k \frac{v_k}{R_k} \\
    i_k  =\frac{v_k}{R_k} 
\end{gather*}
Dunque si è realizzato un sistema che fa la somma pesata dei segnali 
in ingresso 
\begin{gather*}
    v_0 = -i_fR_f = - R_f \sum_k \frac{v_k}{R_k}
\end{gather*}

\section{Integratore}
Questo è un sistema composto da un amplificatore ed un condensatore, 
per cui 
\begin{gather*}
    \frac{Q}{C} = V \ \Longrightarrow \ v_0 = - \frac{Q}{C}
\end{gather*}
In realtà, dato che la carica passa continuamente, 
\begin{gather*}
    v_0 = - \frac{\int_{0}^{t} I \ dt }{C}
\end{gather*}
In questo modo posso integrare il segnale un entrata nell'amplificatore, così che
\begin{gather*}
    v_0 = - \frac{1}{RC}\int_{0}^{t} v_s \ dt 
\end{gather*}

\section{Derivatore}
Un derivatore è un sistema composto da condensatore e resistenza feedback di un 
amplificatore. 
\begin{gather*}
    v_0 = -RC\frac{dv_s}{dt}
\end{gather*} 


\section{Corrente di bias}
La corrente di Bias di un circuito è una corrente 


\section{Tensione di offset}
\begin{gather*}
    \left\{\begin{array}{l}
        v_s -(-v_{off}) = iR_s \\
        v_0 = -v_{off}- iR_s
    \end{array}\right.
\end{gather*}
Allora 
\begin{gather*}
    v_0 = - v_s\frac{R_f}{R_s} - v_{off} \left(\frac{R_s + R_f}{R_s}\right)
\end{gather*}


\end{document}